\textbf{Stage pass rates (Table~\ref{tab:rates}).} Over the runs on frozen code, requirements (stage 1) passed 5 of 6 attempts, implementation (stage 3) 6 of 7, testing (stage 4) 4 of 6, and design (stage 2) only 3 of 9: the design gate was the bottleneck (a file or function that does not exist, a criterion without a row, malformed tool calls while reading the code). Release (stage 5) succeeded every time a run reached it, and CI was green on every pull request. The first-attempt rates are 80\,\%, 40\,\%, 80\,\% and 80\,\%.

\textbf{Hidden tests.} Of the five issues, three reached a pull request. Issue \#4 (None in comparisons) and \#6 (\code{startswith}/\code{endswith}) pass all hidden tests; \#7 (\code{limit}) passes 2 of 4 in round 1 and 3 of 4 after review. The remaining failure is \code{limit=0}, which the specification written by the planner treats as an error and the hidden test treats as an empty result: a reading of the issue, not a coding mistake, and the reviewer approved it against the issue text. Issue \#5 (search results alias the cache) stopped at stage 4: the fix broke existing cache tests in three implementation attempts, the trap the issue contains (an existing test encodes the old behaviour); issue \#8 (vague) stopped at stage 2 after three design attempts. There were no regressions in any pull request.

\textbf{What the reviewer caught, and where it was wrong.} On the pipeline pull requests the reviewer asked for changes twice. On \#7 it found by running its own snippet that the query cache returned before \code{limit} was applied (\code{tinydb/table.py:267}); this is the defect two hidden tests check, and round 2 fixed it. On \#6 it found dead code the tests cannot see (a stray fragment after \code{endswith}), which the hidden tests did not flag. It was wrong, or at least silent, on \#7 round 2: it approved a pull request that still fails one hidden test, because the issue does not say what \code{limit=0} should do. On the seeded pull requests (Table~\ref{tab:e2}) the independent reviewer missed nothing it produced a valid review for, and the invalid reviews are its main failure; early versions (v1, v2) missed the cache-hash bug and produced more invalid reviews, which is why the experiments use only the final version (v3).

\textbf{Prompt injection (\adr{009}).} Testing the screen found that it had been failing open in the recorded runs: its request was refused by our own rate-limit clamp and the gate read the failure as score 0. It is fixed and now fails closed. Measured afterwards, the classifier scores the five real issues at 0.000 to 0.001, blocks two blunt injections (0.993) and misses a polite request, a code-block instruction and a role override (0.001 to 0.018). In an end-to-end test on the fork, two injected issues were stopped at the gate; a third, polite one passed the screen, and the agents treated it as data (the specification quotes it as a source statement, nothing was turned into a requirement, no command ran, no \code{PWNED} reached the diff). The same test showed a git error message writing the job's short-lived authorization header into a stored stop reason; it was fixed (\code{redact}), the comment and files were cleaned, and the old text remains in one runs-branch commit and one run log (\code{experiment/injection\_test.md}).
